% ============================================================================
% Estilo ABNT para o PIM IV — AgroConecta
% Baseado nas regras do manual do PIM (docs/manual.md) e na NBR 14724:2011.
% ============================================================================

% ---------------------------------------------------------------- página ----
\usepackage[a4paper,top=3cm,bottom=2cm,left=3cm,right=2cm]{geometry}

% ------------------------------------------------------------------ fontes ---
\usepackage[T1]{fontenc}
\usepackage[utf8]{inputenc}
\usepackage[portuguese]{babel}
\usepackage{helvet}
\renewcommand{\familydefault}{\sfdefault}

% ------------------------------------------------------------- espaçamento ---
\usepackage{setspace}
\onehalfspacing
\setlength{\parindent}{1.5cm}
\usepackage{indentfirst}
\setlength{\emergencystretch}{1.5em}

% --------------------------------------------------------------- títulos -----
\usepackage{titlesec}
\titleformat{\section}
  {\normalfont\bfseries\fontsize{12}{18}\selectfont}
  {\thesection}{0.8em}{}
\titlespacing*{\section}{0pt}{1.1\baselineskip}{0.7\baselineskip}

\titleformat{\subsection}
  {\normalfont\bfseries\fontsize{12}{18}\selectfont}
  {\thesubsection}{0.8em}{}
\titlespacing*{\subsection}{0pt}{0.9\baselineskip}{0.4\baselineskip}

\titleformat{\subsubsection}
  {\normalfont\itshape\fontsize{12}{18}\selectfont}
  {\thesubsubsection}{0.8em}{}
\titlespacing*{\subsubsection}{0pt}{0.8\baselineskip}{0.3\baselineskip}

% --------------------------------------------------------------- sumário -----
\usepackage{tocloft}
\addto\captionsportuguese{\renewcommand{\contentsname}{SUMÁRIO}}
\renewcommand{\cfttoctitlefont}{\hfill\bfseries\fontsize{12}{18}\selectfont}
\renewcommand{\cftaftertoctitle}{\hfill}
\renewcommand{\cftsecfont}{\normalfont}
\renewcommand{\cftsecpagefont}{\normalfont\bfseries}
\renewcommand{\cftsubsecfont}{\normalfont}
\renewcommand{\cftsubsecpagefont}{\normalfont}
\renewcommand{\cftsecleader}{\cftdotfill{\cftdotsep}}
\setlength{\cftsecindent}{0pt}
\setlength{\cftsubsecindent}{0pt}
\setlength{\cftsubsubsecindent}{0pt}
\setlength{\cftbeforesecskip}{0.4\baselineskip}

% ---------------------------------------------------------------- figuras ----
\usepackage{graphicx}
\usepackage{float}
\usepackage{needspace}
\usepackage{caption}
\captionsetup{
  position=above,
  labelfont=bf,
  labelsep=endash,
  justification=centering,
  singlelinecheck=false,
  font=small,
  skip=6pt
}
\renewcommand{\figurename}{Figura}
\renewcommand{\tablename}{Tabela}

% Fonte de figura/tabela, em 10 pt, logo abaixo do elemento.
\newcommand{\fonte}[1]{%
  \par\vspace{2pt}%
  {\fontsize{10}{12}\selectfont\itshape Fonte: #1\par}%
}

% Atalho: figura com legenda e fonte.
% #1 arquivo, #2 legenda, #3 fonte
\newcommand{\figura}[3]{%
  \begin{figure}[H]
    \centering
    \caption{#2}
    \includegraphics[width=0.92\textwidth]{#1}
    \fonte{#3}
  \end{figure}
}

% Atalho para tabela simples (uma coluna de texto corrida).
% ------------------------------------------------------------- citações ------
% Citação longa: 4 cm da margem esquerda, corpo 10, espaço simples.
\newenvironment{citacao}
  {\begin{spacing}{1}\begin{list}{}{%
     \setlength{\leftmargin}{4cm}\setlength{\rightmargin}{0pt}%
     \setlength{\listparindent}{0pt}\setlength{\itemindent}{0pt}%
     \setlength{\topsep}{0.5\baselineskip}\setlength{\parsep}{0pt}}%
   \item[]\fontsize{10}{12}\selectfont}
  {\end{list}\end{spacing}}

% ------------------------------------------------------------- listas --------
\usepackage{enumitem}
\setlist[itemize]{leftmargin=1.5cm,topsep=0.3\baselineskip,itemsep=0.15\baselineskip}
\setlist[enumerate]{leftmargin=1.5cm,topsep=0.3\baselineskip,itemsep=0.15\baselineskip}

% --------------------------------------------------------------- diagramas ---
\usepackage{tikz}
\usetikzlibrary{positioning, arrows.meta, shapes.geometric, calc, fit, backgrounds}

% --------------------------------------------------------------- tabelas -----
\usepackage{booktabs}
\usepackage{array}
\newcolumntype{L}[1]{>{\raggedright\arraybackslash}p{#1}}
\usepackage{longtable}
\usepackage{multirow}
\setlength{\tabcolsep}{4pt}

% --------------------------------------------------------------- código ------
\usepackage{xcolor}
\usepackage{listings}
\lstdefinestyle{abnt}{
  basicstyle=\ttfamily\fontsize{8}{9.5}\selectfont,
  breaklines=true,
  columns=fullflexible,
  keepspaces=true,
  showstringspaces=false,
  frame=none,
  aboveskip=1.2\baselineskip,
  belowskip=0.8\baselineskip,
  extendedchars=true,
  inputencoding=utf8,
  literate=
    {á}{{\'a}}1 {à}{{\`a}}1 {â}{{\^a}}1 {ã}{{\~a}}1 {ä}{{\"a}}1
    {Á}{{\'A}}1 {À}{{\`A}}1 {Â}{{\^A}}1 {Ã}{{\~A}}1
    {é}{{\'e}}1 {ê}{{\^e}}1 {É}{{\'E}}1 {Ê}{{\^E}}1
    {í}{{\'i}}1 {î}{{\^i}}1 {Í}{{\'I}}1
    {ó}{{\'o}}1 {ô}{{\^o}}1 {õ}{{\~o}}1 {Ó}{{\'O}}1 {Ô}{{\^O}}1 {Õ}{{\~O}}1
    {ú}{{\'u}}1 {ü}{{\"u}}1 {Ú}{{\'U}}1 {Ü}{{\"U}}1
    {ç}{{\c c}}1 {Ç}{{\c C}}1
    {“}{{"}}1 {”}{{"}}1 {‘}{{`}}1 {’}{{'}}1
    {–}{{--}}1 {—}{{---}}1
}
\lstset{style=abnt}

% --------------------------------------------------------------- cabeçalho ---
\usepackage{fancyhdr}
\fancypagestyle{abnt}{%
  \fancyhf{}%
  \fancyhead[R]{\fontsize{10}{12}\selectfont\thepage}%
  \renewcommand{\headrulewidth}{0pt}%
  \renewcommand{\footrulewidth}{0pt}%
}
\fancypagestyle{vazio}{%
  \fancyhf{}%
  \renewcommand{\headrulewidth}{0pt}%
  \renewcommand{\footrulewidth}{0pt}%
}
\fancypagestyle{plain}{%
  \fancyhf{}%
  \renewcommand{\headrulewidth}{0pt}%
  \renewcommand{\footrulewidth}{0pt}%
}

% --------------------------------------------------------------- hyperlinks --
\usepackage[hidelinks,unicode]{hyperref}
\usepackage{xurl}
\hypersetup{
  pdftitle={AgroConecta: plataforma corporativa integrada para gestão da cadeia de suprimentos e inclusão digital na agricultura familiar},
  pdfauthor={[NOME]},
  pdfsubject={Projeto Integrado Multidisciplinar IV}
}

% ------------------------------------------------------- hifenização ---------
\hyphenation{AgroConecta}

% ------------------------------------------------------- comandos auxiliares --
% Título de elemento pré-textual (centralizado, maiúsculo, negrito).
\newcommand{\titulopretextual}[1]{%
  \begin{center}
    \bfseries\fontsize{12}{18}\selectfont #1
  \end{center}
  \vspace{0.5\baselineskip}
}

% Entrada de referência: sem recuo, espaço simples, separada por uma linha.
\newcommand{\refbib}[1]{%
  \par\begingroup
  \setlength{\parindent}{0pt}%
  \setlength{\parskip}{0pt}%
  \noindent #1\par
  \vspace{0.75\baselineskip}%
  \endgroup
}
